% REV_P0 2026-11-24：由 cellcharter_study_conclusion_module_v1_20261122.tex 派生；V1保持不变。
% 按REV_P0结果改写结论1--5及局限；新增数值来源 REV_P0_20261124/derived/T1--T8。
% CellCharter研究专题结论与展望模块（非整篇硕士学位论文总论终稿）
% 对应真实E1/E2,E3c,E3d,E3b,E5a,E5b结果；不预设新算法或统计显著性。
\section{本研究的结论与展望}
\label{sec:cc-conclusion-module-v1}

\subsection{主要实证结论}
本研究在成熟CellCharter空间邻域表示与聚类流程基础上，围绕特征深度、聚类数确定、空间边界以及跨组织评价开展了系统的计算实验。DLPFC部分包含12张切片、3名独立供体；外部数据为6名乳腺癌患者。CellCharter原作者已在同一DLPFC数据上完成基准比较，并对2--4层邻域做过调参\cite{varrone2024,cellcharteranalyses2024}，本研究的25D配置即其最终配置。因此，本研究的贡献是在该基准之上作扩展性的实证评价，不是新方法，也不是对“邻域层数影响ARI”的首次发现。

第一，在固定Oracle-K条件下，每张切片先平均10个技术种子，再对12张切片等权汇总：累计邻域特征从5D增加到25D时，ARI由0.3051提高至0.4418，NMI由0.4133提高至0.5982。其中5D到10D的提升在12/12张切片上方向一致；15D之后，宏平均只略有上升，约半数切片以及供体D3的方向相反。本研究25D的宏平均ARI（0.4418）与原作者官方单切片结果（0.4438）处于同一水平，但两者在重复结构、GMM参数、输入文件和软件版本上存在差异，不构成严格复现。

第二，等维候选B15D（$L_0+L_1+L_4$）是在开发切片151507上看到正向结果后提名的。按预设口径，B相对A15D的12切片平均ARI变化为$+0.0033$；排除两张开发切片后，其余10张切片的平均变化为$-0.0035$（事后汇总）。按供体看，D1为正、D3为负；在外部6名患者的K-blind条件下为$-0.0364$（1/6为正）。B应记为开发集提名、扩展评价中未获稳定支持的负结果。

第三，聚类数设置改变了尺度比较的解释。在E3c预设固定K网格中，每个K下宏平均ARI的尺度排序都是25D $>$ 20D $>$ 15D $>$ 10D，但这种扫描不是无标签模型选择。E3d逐切片AutoK在48个切片--尺度组合中，有10项触及冻结的K=2至10搜索区间边界，选中K与参考类别数一致的情况在每个尺度下只有1/12。本研究每个候选K只做5次重复，少于原作者的10次，且原作者是联合选K。事后评价中，20D的12切片ARI和NMI均值较高（0.4311与0.5942），但供体之间不一致；由于不同尺度实际选中的K不同，这一结果不能作为推荐配置。

第四，空间边界评价呈现与整体聚类指标不同的趋势，以下均为事后分析。严格0-hop F1在5D与10D之间基本持平（0.1441与0.1446），此后下降至25D的0.1179；2-hop容差F1则由0.2863上升至0.5333。严格F1的变化主要来自recall：浅层表示的预测更碎片化，预测边界边数是参考边界的6.4倍（5D），而25D为1.4倍。因此，严格F1较高不等于定位更精确，容差F1较高也不能代替严格定位。这是一项指标定义敏感性观察，不等于已经通过实验确定了“过度平滑”的生物学原因。

第五，与非空间GMM、Leiden及成熟空间方法BANKSY、BayesSpace比较时，CellCharter的表现取决于数据、K策略和指标，不能宣称对所有空间方法普遍占优。本研究冻结的BayesSpace K-blind判据与官方建议的elbow判断不一致\cite{bayesspacevignette}，在DLPFC的36项中有28项、乳腺癌的18项中有18项命中搜索上界；非空间GMM的BIC判据在DLPFC也有32/36项命中上界。这些结果全部保留，但会降低K-blind对比的可比性。此外，对照集未包含原文中单切片表现最好的STAGATE。

第六，在6名乳腺癌患者的外部K-blind评价中，A15D相对5D基础配置的ARI宏平均差约为$+0.0458$、NMI差约为$+0.0583$，两项指标在6/6例患者上方向一致；但在3名ARI低于0.05的患者中，这些差值与种子波动处于同一量级。与BayesSpace比较时，A15D的ARI仅在4/6例患者较高，患者平均NMI差约为$-0.0139$；在Oracle-K下，A15D的ARI低于BayesSpace和Leiden。此结果是将相同研究流程重新用于新组织队列所得的观察，不是将DLPFC已训练模型直接泛化到乳腺癌。

\subsection{研究局限}
本研究在解释层面主要受到以下因素限制：
\begin{enumerate}
  \item DLPFC仅3位供体且存在供体内多个切片，乳腺癌仅6名患者，均不足以支持高精度总体推断。
  \item 25D配置由原作者按参考标签在3张切片上调出，候选B在开发切片上提名，这些切片都不是未接触样本；DLPFC部分参考标注来源及旧开发样本的数据卫生问题需在论文数据来源中保留。
  \item Oracle-K使用参考类别数，不能代表真正未知类别数的无标签应用。
  \item 部分实验中的邻域深度与输入维度同步变化；成熟基线的预处理、计算预算及K-blind规则又存在方法间差异，其中BayesSpace与GMM的似然判据频繁命中上界，对照集未包含STAGATE。
  \item AutoK只做5次稳定性重复，且逐切片运行。
  \item E3d的后验ARI/NMI、E3b的边界与容差分析、排除开发切片的汇总，以及与官方结果的对照，都是在已有模型结果之后追加的事后分析。
  \item 技术种子标准差只反映GMM初始化变异，不是生物学重复。
  \item 病理区域标签与皮层层级标签的语义不同，不能将跨组织ARI绝对值直接对比。
\end{enumerate}

因此，本文展示的是基于已知模型的系统实证评价与可重复的研究流程，而非已验证的新核心聚类算法。在没有更多独立数据和精确因果控制前，不能对边界变化原因、参数最优性或临床应用价值作肯定性普遍结论。

\subsection{后续工作展望}
后续研究可以在以下方向展开：在更多独立供体和组织类型上检验当前观察的适用性；按各方法官方建议的选K方式（例如BayesSpace的elbow判断）以及足够的稳定性重复次数重新设计K-blind对照，并纳入STAGATE等强基线；使用明确的等维和等计算预算对照，进一步区分尺度范围、特征维数和模型随机性；为边界指标建立零模型基准，并结合组织层过渡区域的标注不确定性，避免仅凭单一参考标签和单一指标判断算法表现。

以上属于对当前研究局限的未来研究方向，不表示本研究已经执行了上述扩展。本文当前工作已经完成相应冻结实验及事后分析，毕业论文写作应首先完整呈现已发生的研究事实，而不将计划工作当作已经取得的结果。
